Research on ARC has followed two primary tracks: building stronger solvers and improving evaluation infrastructure. On the solver side, approaches based on program synthesis, search, and large language models continue to advance \citep{alfordNeuralGuidedBidirectionalProgram2022,buttCodeItSelfimprovingLanguage2024,batorskiNSANeurosymbolicARC2025}. On the evaluation side, benchmarks like ConceptARC, MiniARC, and H-ARC reorganize ARC's task space to isolate cognitive concepts, create smaller training settings, and provide robust human baselines \citep{kimPlaygroundsAbstractionReasoning2022,moskvichevConceptARCBenchmarkEvaluating2023,legrisComprehensiveBehavioralDataset2025}. Our proposed environment, JaxARC, supports this second track by providing the simulation speed required to study ARC as a reinforcement learning problem at scale.

Framing ARC as a sequential grid-editing task was first proposed by ARCLE \citep{leeARCLEAbstractionReasoning2025}. By building a Gymnasium-compatible environment, ARCLE showed that PPO agents could learn to solve selected tasks, especially when augmented with model-based priors \citep{leeEnhancingAnalogicalReasoning2024}. However, ARCLE's CPU-bound Python implementation creates a hard limit on experimentation. In a field where progress often depends on large parallel sweeps \citep{luDiscoveredPolicyOptimisation2022,mayorImpactOnPolicyParallelized2025}, a simulator capped at roughly $60,000$ steps per second makes systematic sweeps or meta-learning difficult. JaxARC adopts the conceptual framing of ARCLE but rebuilds the execution layer for hardware accelerators.

This shift mirrors a broader trend where GPU-native simulation has opened new experimental regimes. Continuous control was transformed by Brax \citep{freemanBraxDifferentiablePhysics2021}, board games by Pgx \citep{koyamadaPgxHardwareAcceleratedParallel2023}, and combinatorial optimization by Jumanji \citep{bonnetJumanjiDiverseSuite2023}. In grid-world domains specifically, XLand-MiniGrid \citep{nikulinXLandMiniGridScalableMetaReinforcement2024}, NAVIX \citep{pignatelliNAVIXScalingMiniGrid2025}, and Craftax \citep{matthewsCraftaxLightningFastBenchmark2024} showed that CPU-bound MiniGrid environments artificially limited research scope. NAVIX is most similar to our work: it reworks a widely used Python grid-world suite into a JAX-native implementation and provides a benchmark baseline. JaxARC applies these architectural principles to ARC, though the benchmark introduces unique structural requirements not found in standard grid-worlds. ARC tasks are not fixed MDP instances with a single goal condition; they contain variable-size demonstrations, hidden query targets, exact submit-based scoring, and configurable access to task context. The environment therefore has to preserve benchmark boundaries while still presenting static buffers to the accelerator.

Finally, systematic benchmarking of RL algorithms on challenging domains has proved valuable even when absolute performance is far from human level. The history of deep RL on Atari \citep{mnihPlayingAtariDeep2013} showed that standardized environments and reproducible baselines across games were as important as any individual algorithm advance. Our proposed environment serves a similar role for ARC-style grid editing: it makes the simulator fast enough to systematically measure where current RL methods make progress, where they fail, and how those failures differ across task families.
